Efforts to improve firearm safety through access control have evolved through several distinct technological approaches. 
At the most basic level, traditional solutions rely on purely mechanical safeguards such as trigger locks or integrated 
key-based locking mechanisms. These systems are inexpensive and highly reliable because they do not depend on electronics 
or power, but they do not restrict access to a specific user. Anyone with access to the key or combination can operate 
the firearm, which causes concern toward illegitimate use. In time-sensitive or high-stress situations, the need to manually 
unlock a device with a key can render the safety ineffective in practice, leading many police and military personnel to avoid
using such systems. This feature creates a gap between theoretical safety and real-world usage, which encouraged us to develop 
a system that enforces access control without requiring manual intervention.

Electronic authentication methods represent the next step in this progression. Systems based on PIN entry, external key fobs,
or smartphone connectivity attempt to introduce user-specific control while maintaining relatively simple hardware implementations.
While these approaches improve upon mechanical systems by enabling programmable access, they still require explicit interaction
from the user, such as entering a code or activating a device. This introduces latency and cognitive load that can be unacceptable
in defensive scenarios. Additionally, wireless-based systems may be vulnerable to interception or spoofing, depending on 
their implementation. As a result, these approaches highlight the need for authentication mechanisms that are both secure 
and implicit, allowing the system to verify the user seamlessly without requiring conscious action.

Biometric authentication has emerged as the most prominent alternative in this space, particularly in commercial “smart gun” 
designs. By leveraging physiological traits such as fingerprints, palm vein patterns, or even grip characteristics, biometric 
systems aim to bind firearm operation directly to the identity of the user. This approach is compelling because it eliminates 
the need for external tokens and ensures that only authorized individuals can operate the weapon. However, when 
examined in the context of real-world firearm use, biometrics introduce a number of critical challenges. Sensors such as 
fingerprint scanners are highly sensitive to environmental conditions. Metrics like dirt, sweat, gloves, or moisture can significantly 
degrade performance. Firearms are often used in unpredictable and harsh environments, where such conditions are common, leading
to an increased likelihood of authentication failure.

Latency is another major concern. Biometric systems require time to capture, process, and match data against stored templates. 
Even small delays can be unacceptable in scenarios where immediate response is necessary. Equally important is the issue of 
reliability in terms of false negatives and false positives. A false negative, where an authorized user is denied access, represents
a potentially dangerous failure in a defensive situation. Conversely, a false positive undermines the entire purpose of the
system. These risks are compounded by the long-term degradation of sensors and the variability of human biological signals. In 
addition to these technical concerns, biometric systems tend to be significantly more expensive and power-intensive than 
alternative solutions. Commercial implementations often cost on the order of \$1000 to \$1500, reflecting the complexity of 
most biometric devices. This cost and power burden conflicts directly with the constraints of a small, battery-powered module designed 
to fit within the firing control group of an AR-15.

RFID-based authentication offers a different point in the design space by prioritizing simplicity, robustness, and low power 
consumption. In a passive near-field RFID system, authentication occurs when a tag enters the electromagnetic field generated 
by a reader, allowing the tag to transmit its unique identifier without requiring its own power source. This enables extremely 
low-latency operation and reduces the system's overall energy requirements. RFID systems are also less sensitive to environmental 
contaminants compared to optical or capacitive biometric sensors, making them more reliable in varied conditions. The primary 
tradeoff is that RFID relies on possession rather than inherent identity. A user must carry an external token, such as a ring or 
card, which introduces the possibility of loss or unauthorized use if the token is compromised.

More complex hybrid systems attempt to combine multiple authentication methods, such as pairing RFID with biometrics or behavioral 
sensing, in order to improve security. While these approaches can reduce certain vulnerabilities, they significantly increase system 
complexity, cost, and power consumption. For a compact, retrofittable module operating under strict physical and electrical 
constraints, this hybrid complexity is impractical. As a result, the solution space reflects a series of tradeoffs between security, 
reliability, usability, and feasibility. Biometric systems offer strong identity assurance but struggle with robustness and cost, 
while RFID systems sacrifice some degree of identity binding in exchange for speed, reliability, and practicality. This balance 
makes RFID a practical choice for a system intended to be compact, cost-effective, and deployable in real-world conditions.